\documentclass[10pt]{article}
\usepackage[T1]{fontenc}
\usepackage[utf8]{inputenc}
\usepackage[french]{babel}
\usepackage{amsmath,amssymb,mathrsfs,mathpazo}
\usepackage{geometry}
\geometry{a4paper,margin=2cm}
\usepackage{graphicx,array,tabularx,tabularray,colortbl,multirow,multicol,booktabs,enumitem,pifont,ulem,fancybox,diagbox,cancel}
\usepackage{pstricks,pstricks-add,pst-plot,pst-tree,pst-node,pst-text,pst-func,pst-math,pst-bspline,pst-3dplot,pst-eucl,pst-solides3d}
\let\pscalculate\undefined
\usepackage{pst-bezier}
\usepackage{tikz,pgfplots,tkz-tab}
\usetikzlibrary{arrows,arrows.meta,patterns,calc,positioning,trees,decorations.markings,shapes.geometric,shapes.misc,shapes.arrows,intersections,decorations.pathmorphing,decorations.pathreplacing,backgrounds,shadows,babel}
\pgfplotsset{compat=1.18}
\usepgfplotslibrary{fillbetween}
\usepackage[np]{numprint}
\usepackage{listings,makecell,pas-tableur,twoopt}
\let\barre\undefined
\usepackage[pstricks]{tablvar}
\definecolor{darkgreen}{RGB}{0,100,0}
\providecommand{\red}{\color{red}}
\providecommand{\blue}{\color{blue}}
\providecommand{\green}{\color{green}}
\providecommand{\black}{\color{black}}
\providecommand{\white}{\color{white}}
\providecommand{\gray}{\color{gray}}
\providecommand{\cyan}{\color{cyan}}
\providecommand{\magenta}{\color{magenta}}
\providecommand{\yellow}{\color{yellow}}
\providecommand{\violet}{\color{violet}}
\providecommand{\orange}{\color{orange}}
\providecommand{\euro}{\texteuro}
\providecommand{\EUR}{\texteuro}
\providecommand{\ieme}{\textsuperscript{e}}
\providecommand{\degres}{\ensuremath{^\circ}}
\providecommand{\degree}{\ensuremath{^\circ}}
\providecommand{\vv}[1]{\overrightarrow{#1}}
\providecommand{\Coord}[1]{#1}
\providecommand{\decosix}{}
\providecommand{\starredbullet}{\ensuremath{\bullet}}
\pagestyle{empty}
\setlength{\parindent}{0pt}
\usepackage{adjustbox}\DeclareUnicodeCharacter{2236}{:}\DeclareUnicodeCharacter{2190}{\ensuremath{\leftarrow}}
\providecommand{\R}{\mathbb{R}}
\providecommand{\N}{\mathbb{N}}
\providecommand{\Z}{\mathbb{Z}}
\providecommand{\Q}{\mathbb{Q}}
\providecommand{\D}{\mathbb{D}}
\providecommand{\C}{\mathbb{C}}
\providecommand{\e}{{\mathrm{e}}}
\providecommand{\vect}[1]{\overrightarrow{#1}}
\providecommand{\vectt}[1]{\overrightarrow{\mathrm{#1}}}
\providecommand{\vvt}[1]{\overrightarrow{\mathrm{#1}}}
\providecommand{\np}[2][]{{#2~#1}}
\providecommand{\vv}[1]{{\overrightarrow{#1}}}
\renewcommand{\d}{{\mathrm{d}}}
\providecommand{\degres}{{^\circ}}
\providecommand{\degre}{{^\circ}}
\providecommand{\blue}{\color{blue}}
\providecommand{\red}{\color{red}}
\providecommand{\enskip}{\quad}
\providecommand{\llbracket}{\lbrack\!\lbrack}
\providecommand{\rrbracket}{\rbrack\!\rbrack}
\providecommand{\euro}{{\text{€}}}
\providecommand{\no}{{\text{n°}}}
\providecommand{\interval}[2]{{[#1;#2]}}
\providecommand{\intervFO}[2]{{[#1;#2[}}
\providecommand{\dotfill}{\ldots}

\providecommand{\ds}{\displaystyle}
\providecommand{\textding}[1]{#1}
\providecommand{\Oij}{\(\mathrm{O};\vec{\imath},\vec{\jmath}\)}
\providecommand{\Oijk}{\(\mathrm{O};\vec{\imath},\vec{\jmath},\vec{k}\)}
\providecommand{\Ouv}{\(\mathrm{O};\vec{u},\vec{v}\)}
\providecommand{\intervFF}[2]{\ensuremath{[#1;#2]}}
\renewcommand{\euro}{\ifmmode\text{\texteuro}\else\texteuro\fi}
\DeclareUnicodeCharacter{00B0}{\ifmmode^{\circ}\else\textdegree\fi}
\DeclareUnicodeCharacter{20AC}{\ifmmode\text{\texteuro}\else\texteuro\fi}
% Énoncés : Éducation nationale. Transcription et corrigés : APMEP, Denis Vergès et contributeurs. Découpage et classement : CoopMaths, Sébastien Lozano, Éric Elter et contributeurs. Adaptation de la présentation : Maths Explorer.
% CC BY-SA : https://coopmaths.fr/www/about/
% https://coopmaths.fr/alea/static/bac/2024/tex/bac_2024_06_sujet1_metropole_3.tex
\begin{document}
\noindent\textit{APMEP / CoopMaths. CC BY-SA. Présentation adaptée par Maths Explorer.}\par\medskip
\medskip

L'espace est muni d'un repère orthonormé \Oijk{}.

On considère les points A$(5~;~5~;~0)$, B$(0~;~5~;~0)$, C$(0~;~0~;~10)$ et D$\left(0~;~0 ~;~-\dfrac{5}{2}\right)$.

\begin{center}
	\begin{tikzpicture}[x={(0:6mm)},y={(45:4.6mm)},z={(90:6mm)},>=stealth]
		\fill[gray!20] (5,5,0)--(2.5,2.5,0)--(0,5,0)--(0,0,10)--cycle;
		\draw (5,5,0)node [below right]{A}--
		(0,0,10) node[above left]{C}--
		(0,5,0) node [left]{B}--
		(2.5,2.5,0) node [below] {H}-- cycle (0,0,10)--(2.5,2.5,0);
		\draw[line width = 1.2pt,->] (0,0,0)--(1,0,0) node [below right]{$ \vect{i} $};
		\draw[line width = 1.2pt,->] (0,0,0)--(0,1,0) node [above =2mm]{$ \vect{j} $};
		\draw[line width = 1.2pt,->] (0,0,0)--(0,0,1) node [above left]{$ \vect{k} $};
		\draw (-3,0,0)--(10,0,0) (0,0,-4)--(0,0,11) (0,-4,0)--(0,5,0) (0,9.5,0)--(0,13,0);
		\draw[shift={(0,0,-2.5)}] (-2pt,-2pt)--(2pt,2pt) (2pt,-2pt)--(-2pt,2pt) node [left]{D};
		\draw[shift={(0,0,0)}] (-2pt,-2pt)--(2pt,2pt) (2pt,-2pt)--(-2pt,2pt) node [below right]{O};
		\draw[dashed] (5,5,0)--(0,5,0)--(0,9.5,0);
	\end{tikzpicture}
\end{center}

\begin{enumerate}
	\item \begin{enumerate}
		\item Montrer que $\vect{n_{1}}\begin{pmatrix}1 \\ -1 \\ 0\end{pmatrix}$ est un vecteur normal au plan (CAD).

		\item En déduire que le plan (CAD) a pour équation cartésienne : \quad $x-y=0$.
	\end{enumerate}

	\item On considère la droite $\mathcal{D}$ de représentation paramétrique \renewcommand\arraystretch{1.8}$\left\{\begin{array}{l c l}x&=&\phantom{5 -}\dfrac{5}{2} t \\ y&=&5-\dfrac{5}{2} t \\ z&=&0\end{array}\right.$ où $t \in \R$.\renewcommand\arraystretch{1}


	\begin{enumerate}
		\item On admet que la droite $\mathcal{D}$ et le plan (CAD) sont sécants en un point H. Justifier que les coordonnées de H sont $\left(\dfrac{5}{2}~;~\dfrac{5}{2}~;~0\right)$.

		\item Démontrer que le point H est le projeté orthogonal de B sur le plan (CAD).
	\end{enumerate}

	\item \begin{enumerate}
		\item Démontrer que le triangle ABH est rectangle en H.

		\item En déduire que l'aire du triangle ABH est égale à $\dfrac{25}{4}$.
	\end{enumerate}

	\item \begin{enumerate}
		\item Démontrer que (CO) est la hauteur du tétraèdre ABCH issue de C.

		\item En déduire le volume du tétraèdre ABCH.

		\emph{On rappelle que le volume d'un tétraèdre est donné par : $V=\dfrac{1}{3} \mathcal{B }h$, où $\mathcal{B}$ est l'aire d'une base et $h$ la hauteur relative à cette base.}
	\end{enumerate}

	\item On admet que le triangle ABC est rectangle en B. Déduire des questions précédentes la distance du point H au plan (ABC).
\end{enumerate}

\medskip

\end{document}
